\h*{Reader Roadmap}
\label{reader:roadmap}

This book is organized as a progression from individual instructions to
complete analytical applications. The links below provide five useful entry
routes. If you are new to programming, begin with the first route and continue
in order.

\begin{tcolorbox}[title=Choose a learning route,colframe=orange,colback=white,arc=2mm]
\begin{itemize}
    \item \textbf{Learn Python from the beginning:}
          \booklink{ch:01}{Chapter 1} through \booklink{ch:05}{Chapter 5}.
    \item \textbf{Build reliable, reusable scripts:}
          \booklink{ch:06}{Chapter 6} through \booklink{ch:10}{Chapter 10}.
    \item \textbf{Analyze and communicate data:}
          \booklink{ch:11}{NumPy}, \booklink{ch:12}{pandas}, and
          \booklink{ch:13}{Matplotlib}.
    \item \textbf{Work with geographic data:}
          \booklink{ch:14}{GeoPandas foundations} followed by
          \booklink{ch:15}{geodatasets and applied mapping}.
    \item \textbf{Study computational methods:}
          \booklink{ch:16}{algorithms and efficiency},
          \booklink{ch:17}{numerical optimization}, and
          \booklink{ch:18}{machine-learning foundations}.
\end{itemize}
\end{tcolorbox}

\hh*{Key destinations}

\begin{itemize}
    \item See the \booklink{sec:c13-showcase}{visualization showcase} for a
          compact gallery of analytical plot types.
    \item Follow the \booklink{sec:c13-dashboard}{building-performance
          dashboard} for a longer, practical Matplotlib script.
    \item Use the \booklink{sec:c14-application}{site-access application} to
          connect coordinates, projection, distance, and buffers.
    \item Complete the \booklink{sec:c15-application}{Chicago mapping
          application} to load an educational dataset, align coordinate
          systems, join layers, summarize results, and export a figure.
    \item Compare searches in the
          \booklink{sec:c16-application}{algorithm application}, verify a
          \booklink{sec:c17-application}{bounded optimization model}, and
          evaluate a protected test set in the
          \booklink{sec:c18-application}{machine-learning application}.
\end{itemize}

Every chapter follows the same rhythm: learning outcomes, concepts, syntax
analysis, runnable snippets, visual results where appropriate, guided practice,
solutions, and a checkpoint. Blue text in the PDF is clickable. The footer on
content pages returns to the contents or this roadmap.
